\section{$\operatorname{Lorentz}$群的不可约表示}
在数学上，Poincaré group 与 Lorentz group 的关系是一个标准的半直积结构，表述为：
\begin{equation}
\boxed{
\mathbb{P} = \mathbb{R}^{1,3} \rtimes SO(1,3)
}
\end{equation}
Poincar\'e 变换对应希尔伯特空间中的幺正算符的无穷小形式是：
\begin{equation*}
    {U}(1+\omega,\epsilon)={1}+\frac{i}{2}\omega_{\rho\sigma}\mathscr{J}^{\rho\sigma}-i\epsilon_\rho \mathscr{P}^\rho+\ldots
\end{equation*}
以及Poincar\'e 代数是：
    \begin{equation*}
            \begin{aligned}
        &i\left[\mathscr{J}^{\mu\nu},\mathscr{J}^{\rho\sigma} \right]=
        \eta^{\nu\rho}\mathscr{J}^{\mu\sigma}
        -\eta^{\mu\rho}\mathscr{J}^{\nu\sigma}
        -\eta^{\sigma\mu}\mathscr{J}^{\rho\nu}
        +\eta^{\sigma\nu}\mathscr{J}^{\rho\mu}\\
        &i\left[\mathscr{P}^\mu,\mathscr{J}^{\rho\sigma} \right]=
        \eta^{\mu\rho}\mathscr{P}^\sigma-\eta^{\mu\sigma}\mathscr{P}^\rho\\
        &\left[ \mathscr{P}^\mu,\mathscr{P}^\rho\right]=0
    \end{aligned}
\end{equation*}
去掉半直积的平移部分，便是$Lorentz$变换的无穷小形式以及代数关系是：

\begin{subequations}
    \begin{align}
        M(1+\omega)&={1}+\frac{i}{2}\omega_{\rho\sigma}\mathscr{J}^{\rho\sigma}+\dots\\
         i\left[\mathscr{J}^{\mu\nu},\mathscr{J}^{\rho\sigma} \right]
         &=
        \eta^{\nu\rho}\mathscr{J}^{\mu\sigma}
        -\eta^{\mu\rho}\mathscr{J}^{\nu\sigma}
        -\eta^{\sigma\mu}\mathscr{J}^{\rho\nu}
        +\eta^{\sigma\nu}\mathscr{J}^{\rho\mu}
    \end{align}
\end{subequations}
$Lorentz$代数的三维角动量和$boost$\footnote{别忘了$boost$是不守恒且反厄米的}是：
\begin{subequations}
    \begin{align}
        &\mathbf{J} = \{\mathscr{J}^{23}, \mathscr{J}^{31}, \mathscr{J}^{12}\}\\
        &\mathbf{K} = \{\mathscr{J}^{01}, \mathscr{J}^{02}, \mathscr{J}^{03}\}
    \end{align}
\end{subequations}
相应的对易关系是：
\begin{subequations}
    \begin{align}
        [J_{i}, J_{j}] &= i\epsilon_{ijk} J_{k}  \\
        [J_{i}, K_{j}] &= i\epsilon_{ijk} K_{k} \\
        [K_{i}, K_{j}] &= -i\epsilon_{ijk} J_{k} 
    \end{align}
\end{subequations}
\text{\color{blue}利用线性组合，$Lorentz$代数变为}：
\begin{subequations}
    \begin{align}
        &A_i=\frac{1}{2}(J_i-iK_i)\\
        &B_i=\frac{1}{2}(J_i+iK_i)
    \end{align}
\end{subequations}
这俩个算符的对易关系是：
\begin{subequations}
    \begin{align}
        &[A_i,A_j]=i\epsilon_{ijk} A_{k}\\
        &[B_i,B_j]=i\epsilon_{ijk} B_{k}\\
        &[A_i,B_j]=0
    \end{align}
\end{subequations}
这时，\color{blue}$Lorentz$群可以被视为两个“角动量群”。
值得注意：给“角动量群”加上引号，是因为生成元
并非厄米的，因此它们生成的群严格来说不是通常的幺正SU(2)群，而是其复化版本。的直积\normalcolor；因此复化后的$Lorentz$代数分解为
\begin{equation}
\boxed{\mathfrak{so}(1,3)_{\mathbb C}
\simeq
\mathfrak{su}(2)\oplus\mathfrak{su}(2)}
\end{equation}
将完整的HLG代数分解为两个对易子群，有效地解决了$Lorentz$群所有有限维表示的分类问题。HLG的一般表示可以用自旋对$(A,B)$标记，它们的取值是：
\begin{subequations}
    \begin{align}
        &a=-A,-A+1,\dots,+A\\
        &b=-B,-B+1,\dots,+B
    \end{align}
\end{subequations}
两组角动量的基底为：
\begin{equation}
    \ket{A,a}\oplus\Ket{B,b}=\Ket{A,B;a,b}\overset{\text{简写为}}{=}\Ket{a,b}
\end{equation}
\begin{subequations}
\begin{align}
    \braket{a',b'|\mathbf{A}|a,b}
    &= 
    \braket{a'|\mathbf{J}^{A}|a}\,\delta_{bb'} ,\\
    &\to J^A_3\ket{a}
    = a\ket{a},\\
    &\to (J^A_1 \pm iJ^A_2)\ket{a}
    = 
    \sqrt{(A \mp a)(A \pm a + 1)}\,\ket{a \pm 1}
\end{align}
\begin{align}
    \braket{a',b'|\mathbf{B}|a,b}
    &=
    \delta_{aa'}\,\braket{b'|\mathbf{J}^{B}|b},\\
    &\to J^B_3\ket{b}
    = b\ket{b},\\
    &\to (J^B_1 \pm iJ^B_2)\ket{b}
    = 
    \sqrt{(B \mp b)(B \pm b + 1)}\,\ket{b \pm 1}.
\end{align}
\end{subequations}
\text{\color{blue}所以场和角动量耦合理论息息相关}：一般的总角动量$\boldsymbol{J}$可以通过$\boldsymbol{J}=\boldsymbol{J}^A+\boldsymbol{J}^B$,取值遍及两个角动量差到和：
\begin{equation}
    |A-B|\leq j\leq |A+B|
\end{equation}
这时，将(A,B)翻译成各种物理场;给出表格
\begin{table}[htbp]
    \centering
    \renewcommand{\arraystretch}{1.4}
    % 如果你需要图表标题是中文可以自己把 caption 改回来
    \caption{Representations of the Lorentz Group and Corresponding Physical Fields}
    \vspace{0.2cm}
    \begin{tabular}{c c l}
        \toprule
        \textbf{$(A,B)$ Representation} & \textbf{Max Spin $j$} & \textbf{Field Type} \\
        \midrule
        $(0, 0)$ & $0$ & Scalar Field \\
        
        $(\frac{1}{2}, 0)$ or $(0, \frac{1}{2})$ & $\frac{1}{2}$ & Left/Right-handed Weyl Spinor \\
        
        $(\frac{1}{2}, 0) \oplus (0, \frac{1}{2})$ & $\frac{1}{2}$ & Dirac Spinor or  Majorana Spinor \\
        
        $(\frac{1}{2}, \frac{1}{2})$ & $1$ & Vector Field \\
        
        $(1, 0)$ or $(0, 1)$ & $1$ & (Anti-)Self-dual Tensor Field \\
        
        $(1, 0) \oplus (0, 1)$ & $1$ & Antisymmetric Tensor Field (2-form) \\
        
        $(1, \frac{1}{2}) \oplus (\frac{1}{2}, 1)$ & $\frac{3}{2}$ & Rarita-Schwinger Field (Vector-Spinor) \\
        
        $(1, 1)$ & $2$ & Symmetric Traceless Tensor Field \\
        \midrule
        $(A, A)$ & $2A$ & Totally Symmetric Traceless Tensor Field of Rank $2A$ \\
        \bottomrule
    \end{tabular}
    \label{tab:fields_representation_english}
\end{table}
\begin{comment}
    begin{center}
    各类物理场的自由度
\end{center}
\vspace{0.2cm}
在量子场论中，场在$Lorentz$群表示下的“Off-shell”数学分量，与实际满足运动方程的“On-shell”物理传播自由度是截然不同的。特别是对于无质量规范场，规范对称性会消除非物理的纵向极化。具体讨论如下：

\begin{itemize}
    \setlength{\itemsep}{0.8em} % 增加条目间的垂直间距，用留白代替加粗来区分层次

    \item 标量场 $(0,0)$，自旋 $j=0$。
    \begin{itemize}
        \item 物理上对应无方向性的标量粒子，如 Klein-Gordon 场、Higgs Boson。
        \item 离壳维度为 $(2\times0+1)^2=1$，在壳（无论有无质量）始终保持 $1$ 个物理自由度。
    \end{itemize}

    \item Weyl 旋量场 $(\frac{1}{2}, 0)$ 或 $(0, \frac{1}{2})$，自旋 $j=1/2$。
    \begin{itemize}
        \item 描述具有确定手征性（左手或右手）的无质量费米子，是构建标准模型的基本模块。
        \item 具有 $2$ 个物理自由度（对应螺旋度 $\pm 1/2$）。
    \end{itemize}

    \item Dirac / Majorana 旋量场 $(\frac{1}{2}, 0) \oplus (0, \frac{1}{2})$，自旋 $j=1/2$。
    \begin{itemize}
        \item Dirac 场：由左右手外尔旋量直和而成，描述有质量的费米子（如电子、夸克）。离壳具有 $2+2=4$ 个复分量；在壳具有 $4$ 个物理自由度（正/反粒子 $\times$ 自旋上下）。
        \item Majorana 场：对狄拉克场施加实数条件（要求反粒子就是其自身，如中微子假说），自由度减半，仅剩 $2$ 个物理自由度。
    \end{itemize}

    \item 四维矢量场 $(\frac{1}{2}, \frac{1}{2})$，自旋 $j=1$。
    \begin{itemize}
        \item 对应四维协变矢量 $A_\mu$。离壳分量数为 $(2\times\frac{1}{2}+1)^2 = 4$。
        \item 若有质量（ W/Z 玻色子）：存在纵向极化，物理状态数为 $2j+1 = 3$，即 $3$ 个物理自由度。
        \item 若无质量（光子\footnote{光子到底装进无质量矢量场还是反对称张量场，或许是一个值得讨论的问题
        }）：规范对称性消除纵极化，仅剩横向极化，具有 $2$ 个物理自由度。
    \end{itemize}

    \item 反对称张量场 $(1, 0) \oplus (0, 1)$，自旋 $j=1$ 的反对称部分。
    \begin{itemize}
        \item 对应二阶反对称张量（2-form）$B_{\mu\nu} = -B_{\nu\mu}$，如弦论中的 Kalb-Ramond 场。离壳有 $3+3=6$ 个独立分量。
        \item 若有质量：具有 $3$ 个物理自由度。
        \item 若无质量：它只有 $1$ 个物理自由度。这可通过两种视角理解：
        \begin{enumerate}
            \item 小群横向截面法：无质量粒子的极化只能存在于垂直运动方向的 $2\times 2$ 横向平面内。在此平面内的反对称矩阵只有唯一的独立分量（$B_{12} = -B_{21}$）。
            \item Hodge Duality：在四维时空中，2-形式场 $B_{\mu\nu}$ 的场强是 3-形式。通过 $\star$ 对偶操作，3-形式等价于 1-形式的导数。因此它在物理上完全等同于一个无质量标量场。
        \end{enumerate}
    \end{itemize}

    \item Rarita-Schwinger 场 $(1, \frac{1}{2}) \oplus (\frac{1}{2}, 1)$，自旋 $j=3/2$。
    \begin{itemize}
        \item 描述带有矢量指标和旋量指标的场 $\Psi_\mu^\alpha$，是超对称理论中引力微子（Gravitino）的标准描述。离壳分量数为 $3\times2 + 2\times3 = 12$。
        \item 若有质量，具有 $4$ 个物理自由度；若无质量，具有 $2$ 个物理自由度。
    \end{itemize}

    \item 二阶对称无迹张量场 $(1, 1)$，自旋 $j=2$。
    \begin{itemize}
        \item 对应度规微扰 $h_{\mu\nu}$，描述引力场。
        \item 离壳计数：一个一般的四维二阶对称张量有 $4 \times 5 / 2 = 10$ 个分量。扣除 $1$ 个无迹条件（$h^\mu_\mu = 0$）后，剩余 $9$ 个分量，完美吻合群论维度 $(2\times 1+1)^2 = 9$。
        \item 若有质量（Fierz-Pauli 场）：具有 $2j+1 = 5$，即 $5$ 个物理自由度。
        \item 若无质量（广义相对论的引力波）：微分同胚不变性去除了多余极化，仅保留加号（$+$）和叉号（$\times$）极化，具有 $2$ 个物理自由度。
    \end{itemize}

    \item 完全对称无迹张量场 $(A, A)$，广义高自旋场，自旋 $j=2A$（整数）。
    \begin{itemize}
        \item 张量构造与无迹操作的本质：这是一个带有 $2A$ 个$Lorentz$指标的完全对称张量 $\Phi_{\mu_1 \mu_2 \dots \mu_{2A}}$（如 Fronsdal 场）。为了使其构成$Lorentz$群的不可约表示，必须扣除所有包含低阶表示的“迹”。即要求对任意两个指标缩并都为零：$\eta^{\mu_1 \mu_2} \Phi_{\mu_1 \mu_2 \dots \mu_{2A}} = 0$。
        \item 离壳自由度完美闭环计算：
        \begin{itemize}
            \item 在 4 维时空中，秩为 $N=2A$ 的完全对称张量有 $\binom{4+N-1}{N}$ 个分量。
            \item 它的迹本身是一个秩为 $N-2$ 的完全对称张量，拥有 $\binom{4+(N-2)-1}{N-2}$ 个分量。
            \item 无迹操作相当于强行减去这些分量：$\binom{N+3}{N} - \binom{N+1}{N-2} = (N+1)^2$。
            \item 代入 $N=2A$，得到离壳维度正好是 $(2A+1)^2$，与$Lorentz$群表示 $(A,A)$ 的维度公式严丝合缝。
        \end{itemize}
        \item 若有质量：具有 $2(2A)+1 = 4A+1$ 个物理自由度。
        \item 若无质量：所有自旋 $j \ge 1$ 的无质量规范场都只有两个螺旋度状态 $\lambda = \pm 2A$，因此恒为 $2$ 个物理自由度。
    \end{itemize}\end{itemize}
\end{comment}


\begin{center}
    \vspace{2ex}
    \text{$Lorentz$群的表示：非紧致性与空间反演}\\
    \vspace{2ex}
\end{center}
在通过 $SU(2)_L \times SU(2)_R$ 代数构建$Lorentz$群的有限维表示 $(A,B)$ 后，有两个深刻的数学物理性质决定了量子场论的具体形态。

\paragraph{1. 非紧致性与有限维表示的非幺正性}

$Lorentz$群 $SO(1,3)$ 与四维欧几里得旋转群 $SO(4)$ 有着本质的拓扑差异。$SO(4)$ 的群参数空间是闭合有界的（紧致群），而 $SO(1,3)$ 包含了 Boost 变换——其快度参数
\begin{equation}
    \phi = \text{arctanh}\frac{v}{c} = \frac{1}{2}\ln\frac{1+v/c}{1-v/c}
\end{equation}
在 $v\to c$ 时趋于无穷大。这意味着 $SO(1,3)$ 的参数空间是无界的，因此它是一个\textbf{非紧致群}。

根据群表示论中的基本定理\footnote{对于非紧致的半单李群，所有忠实的有限维表示都是非幺正的。这是紧致群才拥有的 Peter-Weyl 定理不再成立的直接后果。}，非紧致的半单李群\textbf{不存在有限维的幺正表示}（平凡表示 $(0,0)$ 除外）。

这一抽象定理在 $(A,B)$ 表示中有极其具体的体现。回顾生成元的定义：
\begin{equation}
    \mathbf{A} = \frac{1}{2}(\mathbf{J}+i\mathbf{K}),\quad \mathbf{B} = \frac{1}{2}(\mathbf{J}-i\mathbf{K})
\end{equation}
其中 $\mathbf{J}$ 是角动量（旋转生成元），$\mathbf{K}$ 是 Boost 生成元。由于 $\mathbf{J}$ 厄米而 $\mathbf{K}$ 必须反厄米（这正是非紧致性的体现），我们有：
\begin{itemize}
    \item 旋转生成元 $\mathbf{J} = \mathbf{A}+\mathbf{B}$ 是\textbf{厄米}的 $\Rightarrow$ 旋转矩阵 $e^{i\boldsymbol{\theta}\cdot\mathbf{J}}$ 是幺正的。
    \item Boost 生成元 $\mathbf{K} = -i(\mathbf{A}-\mathbf{B})$ 是\textbf{反厄米}的 $\Rightarrow$ Boost 矩阵 $e^{-i\boldsymbol{\phi}\cdot\mathbf{K}} = e^{-\boldsymbol{\phi}\cdot(\mathbf{A}-\mathbf{B})}$ \textbf{不是幺正的}。
\end{itemize}

这意味着，我们在表 5.1 中列出的所有有限维物理场（$2\times 2$ 的旋量、$4\times 4$ 的矢量等），其$Lorentz$变换矩阵在 Boost 下都是非幺正的。

但这并不构成物理上的灾难。原因在于：我们关注的对象是\textbf{场算符}，而不是\textbf{态矢量}。态矢量生活在无穷维的 Hilbert 空间中，其$Lorentz$变换由 Wigner 的诱导表示给出，\textbf{是}幺正的。

换言之：\textbf{粒子态的幺正性由 Poincaré 群的无穷维表示保障；场算符的非幺正性则是有限维表示为洛伦兹协变性所付出的代价。}

\paragraph{2. 空间反演与直和表示}

如果理论要求包含空间反演（宇称），则必须存在一个表示矩阵 $\beta$ 使得生成元在其作用下满足：
\begin{equation}
    \beta\,\mathbf{J}\,\beta^{-1} = +\mathbf{J},\quad \beta\,\mathbf{K}\,\beta^{-1} = -\mathbf{K}
\end{equation}
物理上这很容易理解：空间反演下坐标 $\mathbf{x}\to-\mathbf{x}$，动量 $\mathbf{p}\to-\mathbf{p}$，因此角动量 $\mathbf{J}=\mathbf{r}\times\mathbf{p}\to+\mathbf{J}$（赝矢量不变），而 Boost 生成元 $\mathbf{K}\sim t\mathbf{p}-\mathbf{x}E\to-\mathbf{K}$（极矢量反号）。

$\mathbf{K}$ 的反号使得左右手生成元在宇称下互相交换：
\begin{equation}
    \mathcal{P}:\quad \mathbf{A}=\frac{1}{2}(\mathbf{J}+i\mathbf{K}) \longleftrightarrow \frac{1}{2}(\mathbf{J}-i\mathbf{K})=\mathbf{B}
\end{equation}
这意味着空间反演将 $(A,B)$ 表示直接映射为 $(B,A)$ 表示：
\begin{equation}
    \mathcal{P}:\quad (A,B) \longrightarrow (B,A)
\end{equation}
因此我们立即得出结论：
\begin{itemize}
    \item 若 $A=B$（如标量场 $(0,0)$、矢量场 $(1/2,1/2)$、引力场 $(1,1)$），表示在宇称下天然封闭。
    \item 若 $A\neq B$（如 Weyl 旋量 $(1/2,0)$ 或 $(0,1/2)$），单一不可约表示\textbf{无法}提供宇称对称性。
\end{itemize}

这正是为什么自然界中的手征费米子（Weyl 旋量）破坏了宇称对称性——弱相互作用中的左手流 $\bar{\psi}_L\gamma^\mu\psi_L$ 只涉及 $(1/2,0)$ 表示，不包含 $(0,1/2)$。

为了在理论中恢复空间反演对称性（如 QED 和 QCD），我们必须将不等价的表示进行直和，构造出 $(A,B)\oplus(B,A)$ 的形式。典型的例子包括：
\begin{align}
    \text{Dirac 场：}\quad &\left(\frac{1}{2},0\right)\oplus\left(0,\frac{1}{2}\right) \\
    \text{电磁场强张量：}\quad &(1,0)\oplus(0,1)
\end{align}
这正是狄拉克旋量必须是四分量的深层含义，也是电磁场张量 $F_{\mu\nu}$ 必须同时包含电场 $\mathbf{E}$（极矢量）和磁场 $\mathbf{B}$（赝矢量）的根本原因。
\begin{center}
    \vspace{2ex}
无质量粒子的 $(A,B)$ 表示分析
\vspace{2ex}

\end{center}

在 5.2 节的前半部分，我们通过 $SU(2)_L \times SU(2)_R$ 代数对 Lorentz 群的有限维表示进行了完整分类，并列出了各类 $(A,B)$ 表示对应的物理场。对于有质量粒子，表示空间 $V^{(A)} \otimes V^{(B)}$ 中的全部 $(2A+1)(2B+1)$ 个分量都参与物理——静止系中的系数由角动量加法（Clebsch--Gordan 系数）决定，粒子具有 $2j+1$ 个自旋投影。

对于无质量粒子，情况发生了根本变化。第 3 章中我们已经确定无质量粒子的 Wigner 小群为 $ISO(2)$，并通过排除连续自旋表示将物理表示约化为 $SO(2)$ 的一维表示——螺旋度 $\lambda$。现在要在 $(A,B)$ 表示的框架内，回答一个关键的表示论问题：\textbf{$(A,B)$ 表示中的哪些分量能够与螺旋度 $\lambda$ 的无质量粒子耦合？}


\subparagraph{$ISO(2)$ 生成元在 $(A,B)$ 表示中的分解}

回忆 $ISO(2)$ 小群（标准动量 $k^\mu = (\kappa,0,0,\kappa)$）的三个生成元为
\[
J_3\,,\qquad \Pi_1 = K_1 - J_2\,,\qquad \Pi_2 = K_2 + J_1
\]
满足李代数 (3.32)：
\[
[J_3, \Pi_1] = i\Pi_2\,,\qquad [J_3, \Pi_2] = -i\Pi_1\,,\qquad [\Pi_1, \Pi_2] = 0
\]

在 $(A,B)$ 表示中，Lorentz 群的生成元分解为（参见 (5.27)(5.55)）
\begin{equation}
\mathbf{J} = \mathbf{J}^{(A)} + \mathbf{J}^{(B)}\,,\qquad
\mathbf{K} = -i\mathbf{J}^{(A)} + i\mathbf{J}^{(B)}
\end{equation}
将上式代入 $\Pi_1 = K_1 - J_2$ 和 $\Pi_2 = K_2 + J_1$ 的定义，引入升降算符 $J^{(X)}_\pm = J^{(X)}_1 \pm iJ^{(X)}_2$，构造组合 $\Pi_1 \pm i\Pi_2$。

\textbf{计算 $\Pi_1 + i\Pi_2$：}
\begin{align}
\Pi_1 + i\Pi_2 &= (K_1 + iK_2) + (iJ_1 - J_2) \notag \\
&= \bigl(-iJ_1^{(A)} + J_2^{(A)} + iJ_1^{(B)} - J_2^{(B)}\bigr)
+ \bigl(iJ_1^{(A)} + iJ_1^{(B)} - J_2^{(A)} - J_2^{(B)}\bigr) \notag \\
&= 2iJ_1^{(B)} - 2J_2^{(B)} = 2i\bigl(J_1^{(B)} + iJ_2^{(B)}\bigr)
= 2i\,J_+^{(B)}
\end{align}

\textbf{计算 $\Pi_1 - i\Pi_2$：}
\begin{align}
\Pi_1 - i\Pi_2 &= (K_1 - iK_2) + (-iJ_1 - J_2) \notag \\
&= \bigl(-iJ_1^{(A)} - J_2^{(A)} + iJ_1^{(B)} + J_2^{(B)}\bigr)
+ \bigl(-iJ_1^{(A)} - iJ_1^{(B)} - J_2^{(A)} - J_2^{(B)}\bigr) \notag \\
&= -2iJ_1^{(A)} - 2J_2^{(A)} = -2i\bigl(J_1^{(A)} - iJ_2^{(A)}\bigr)
= -2i\,J_-^{(A)}
\end{align}

总结：
\begin{equation}
\boxed{
\Pi_1 + i\Pi_2 = 2i\,J_+^{(B)}\,,\qquad
\Pi_1 - i\Pi_2 = -2i\,J_-^{(A)}
}
\end{equation}
$ISO(2)$ 的两个``平移''生成元在 $(A,B)$ 表示中\textbf{完全分解}为 $J_+^{(B)}$（自旋 $B$ 的升算符）和 $J_-^{(A)}$（自旋 $A$ 的降算符）。

\subparagraph{旋转条件与``平移''条件对 $(A,B)$ 分量的约束}

正如有质量粒子的静止系系数 $u_{ab}(0,\sigma)$ 由小群 $SU(2)$ 的旋转条件决定（交织方程 (5.45) 最终给出 CG 系数），无质量粒子标准动量处的系数 $u_{ab}(\mathbf{k},\lambda)$ 由小群 $ISO(2)$ 的\textbf{两类}条件决定：

\textbf{1. 旋转条件（$J_3$ 部分）：} $J_3$ 在 $(A,B)$ 表示中是对角的，$J_3 = J_3^{(A)} + J_3^{(B)}$，本征值为 $a+b$。要求场的系数在绕 $z$ 轴旋转 $\theta$ 后获得相位 $e^{i\lambda\theta}$（与粒子态的螺旋度变换一致），得到
\begin{equation}
u_{ab}(\mathbf{k},\lambda) = 0 \quad \text{除非} \quad \lambda = a + b
\end{equation}
这与有质量情况中的旋转条件 $\sigma = a + b$（方程 (5.45a)）形式完全相同。对于有质量粒子，$\sigma$ 可以取 $|A-B|$ 到 $A+B$ 之间的任何值；差异出现在下面的``平移''条件中。

\textbf{2. ``平移''条件（$\Pi_1, \Pi_2$ 部分）：} 在 3.4 节中我们已经论证，物理表示要求 $\Pi_1 = \Pi_2 = 0$（排除连续自旋）。利用上面的分解结果，$\Pi_1\,u = 0$ 且 $\Pi_2\,u = 0$ 等价于
\begin{equation}
J_-^{(A)}\,u = 0\,,\qquad J_+^{(B)}\,u = 0
\end{equation}

$J_-^{(A)}$ 是自旋 $A$ 表示中的降算符。使降算符给出零的\textbf{唯一}本征态是最低权态 $a = -A$：
\[
J_-^{(A)}|A,-A\rangle = 0
\]
因此 $u_{\bar{a}b} = 0$ 除非 $a = -A$。

$J_+^{(B)}$ 是自旋 $B$ 表示中的升算符。使升算符给出零的唯一本征态是最高权态 $b = +B$：
\[
J_+^{(B)}|B,+B\rangle = 0
\]
因此 $u_{a\bar{b}} = 0$ 除非 $b = +B$。

综合旋转条件 $\lambda = a+b$ 与``平移''条件 $a = -A$、$b = +B$：
\begin{equation}
\boxed{
u_{ab}(\mathbf{k},\lambda) = 0 \quad \text{除非}\quad a = -A,\;\; b = +B\,,\quad \text{从而}\quad \lambda = B - A
}
\end{equation}

对 $v$ 系数完全类似的分析（$v$ 方程中螺旋度相位取复共轭）给出 $v_{ab}(\mathbf{k},\lambda) = 0$ 除非 $a = -A$、$b = +B$，且 $-\lambda = a+b = B-A$，即 $v$ 部分对应\textbf{螺旋度 $\lambda = A-B = -(B-A)$} 的反粒子。

\subparagraph{核心结论}

\begin{center}
\fbox{\parbox{0.88\textwidth}{
\textbf{$(A,B)$ 表示与螺旋度的对应关系}

一个按照 Lorentz 群的 $(A,B)$ 不可约表示变换的场：
\begin{itemize}
\item 湮灭螺旋度 $\lambda = B - A$ 的粒子；
\item 产生螺旋度 $\lambda = A - B = -(B-A)$ 的反粒子。
\end{itemize}
\vspace{4pt}
整个 $(2A+1)(2B+1)$ 维表示空间中，仅有\textbf{一个分量} $u_{-A,+B}$ 存活。
}}
\end{center}

\subparagraph{与有质量情况的对比}

\begin{center}
\begin{tabular}{c|c|c}
\hline
& 有质量粒子 & 无质量粒子 \\
\hline
小群 & $SU(2)$ & $ISO(2)\xrightarrow{\Pi=0}SO(2)$ \\[2pt]
标准动量系数 & $u_{ab}(0,\sigma) = \dfrac{1}{\sqrt{2m}}\langle A,a;B,b|j,\sigma\rangle$
& $u_{ab}(\mathbf{k},\lambda) \propto \delta_{a,-A}\,\delta_{b,+B}$ \\[6pt]
非零分量数 & $(2A+1)(2B+1)$（全部） & $1$（仅 $a=-A,\,b=+B$） \\[2pt]
内部自由度 & $\sigma = -j,\ldots,+j$（共 $2j+1$ 个） & $\lambda = B-A$（仅 $1$ 个） \\[2pt]
决定系数的机制 & 角动量加法（CG 系数） & 升降算符的零化条件 \\
\hline
\end{tabular}
\end{center}

有质量粒子具有静止系，$SU(2)$ 小群的全部 $2j+1$ 维空间都被利用，系数由 Clebsch--Gordan 系数给出。无质量粒子没有静止系，$ISO(2)$ 小群的``平移''部分（$\Pi_1,\Pi_2$）施加了额外约束，将 $(2A+1)(2B+1)$ 维空间\textbf{冻结}到仅一个分量——这就是无质量粒子只有两个螺旋度 $\pm j$（而非 $2j+1$ 个极化）的\textbf{群论根源}。

\subparagraph{各类 $(A,B)$ 表示允许的螺旋度}

利用 $\lambda = B - A$，列出常见表示能描述的无质量粒子：

\begin{center}
\begin{tabular}{c|c|c|l}
\hline
$(A,B)$ & 场的类型 & $\lambda = B-A$ & 物理对应 \\
\hline
$(0,0)$ & 标量 & $0$ & 无质量标量 \\
$(\frac{1}{2},0)$ & 左手 Weyl 旋量 & $-\frac{1}{2}$ & 左手中微子 \\
$(0,\frac{1}{2})$ & 右手 Weyl 旋量 & $+\frac{1}{2}$ & 右手反中微子 \\
$(\frac{1}{2},\frac{1}{2})$ & 四矢量 & $0$ & \textbf{不能}描述光子 \\
$(1,0)$ & 自对偶张量 & $-1$ & 螺旋度 $-1$ \\
$(0,1)$ & 反自对偶张量 & $+1$ & 螺旋度 $+1$ \\
$(2,0)$ & --- & $-2$ & 螺旋度 $-2$（引力子） \\
$(0,2)$ & --- & $+2$ & 螺旋度 $+2$（引力子） \\
\hline
\end{tabular}
\end{center}

特别地：
\begin{itemize}
\item 描述螺旋度 $\pm j$ 的自共轭无质量粒子需要直和表示 $(j,0)\oplus(0,j)$。

\item 四矢量 $(\frac{1}{2},\frac{1}{2})$ 只能描述螺旋度零——这预示着构造螺旋度 $\pm 1$ 的无质量矢量场时将不可避免地遇到困难，并最终导致规范不变性的引入。这些内容将在 5.4 节构造零质量粒子场时详细展开。

\item 自旋 $j$ 的无质量粒子的 $(A,A+j)$ 场是 $(0,j)$ 场的 $2A$ 阶导数，$(B+j,B)$ 场是 $(j,0)$ 场的 $2B$ 阶导数。因此本质独立的最简单场只有 $(j,0)$ 和 $(0,j)$。
\end{itemize}
